\section{Related Work}

\paragraph{Action chunking policies.}
Chunked action prediction is standard in modern imitation learning. ACT
predicts chunks with a transformer and blends them with temporal ensembling
\citep{zhao2023act}. Diffusion Policy denoises an action sequence and executes
a prefix before replanning \citep{chi2023diffusion}. Generalist
vision-language-action models such as GR00T N1 produce 16-step chunks from a
flow-matching head \citep{gr00t2025}. In all of these the executed chunk
length is a fixed hyperparameter, chosen once per task.

\paragraph{Compounding error and the theory of chunking.}
The reason short chunks help at all is the compounding of error under a policy
that acts without observing, analyzed for imitation learning by
\citet{ross2010efficient} and \citet{ross2011dagger}, whose bound grows
quadratically in the horizon in the worst case. In continuous action spaces the
situation is sharper: \citet{simchowitz2025pitfalls} show that even with
exponentially stable dynamics and a smooth deterministic expert, any smooth
imitator can suffer execution error exponential in the horizon relative to its
error on the training distribution. Closest to this paper,
\citet{zhang2025chunking} prove that action chunking is a remedy, because a
chunk longer than a threshold set by the contraction rate restores exponential
stability of the learned closed loop and makes the trajectory error
horizon-free. Their analysis and ours meet at the same inequality and pull in
opposite directions on the same knob. Theirs is a lower bound on chunk length,
derived under a single contraction rate for the system; ours is the upper bound
that the same inequality implies when the local rate is positive, and our
measurement is the state-dependent counterpart of their constant. Section
\ref{sec:method} makes the connection explicit, and Claim 1 is the empirical
observation that no single rate describes an episode.

\paragraph{Adaptive execution horizons.}
A recent line of work varies the executed horizon at test time, and the
distinguishing feature is the signal each one switches on. Bidirectional
decoding re-samples chunks and selects for closed-loop coherence
\citep{liu2024bidirectional}. DEHP trains an execution-horizon branch with
chunk-level reinforcement learning over a frozen base policy
\citep{zhao2026dehp}. AutoHorizon reads the policy's action self-attention as a
confidence signal \citep{wang2026autohorizon}. PACE cuts chunks at low-speed
valleys of the predicted trajectory \citep{nie2026pace}. DVAC uses the variance
of the clean-action estimate along the denoising trajectory of a flow policy
\citep{feng2026dvac}, and SGAC compares the queued chunk against a freshly
sampled one \citep{so2025sgac}. HiPolicy predicts chunks at several frequencies
and gates them by entropy \citep{zhang2026hipolicy}, and \citet{wen2026adaptive}
learn chunk length end-to-end from visual context.

These methods share our motivation that manipulation is stage-structured, and
PACE states that intuition explicitly. The difference is what gets measured.
Confidence, attention, entropy and sample agreement describe the policy's
uncertainty about its own output. Speed valleys describe the kinematics of the
predicted path. None of them measure how fast the plant amplifies an action
error, which is the quantity chunk length actually trades against, and none
produce per-timestep labels that exist independently of the controller that
consumes them. We measure these proxies against $\lambda$ at the same states in
Section~\ref{sec:proxies}.

\paragraph{Stability analysis in control and learning.}
Finite-time Lyapunov exponents are a classical tool for diagnosing trajectory
divergence and coherent structure in flows \citep{haller2015lcs}, and have been
used to analyze the sensitivity of MPC and reinforcement learning policies
\citep{krishna2023ftle}. We bring the tool to a different question: not whether
a controller is stable overall, but which chunk length a chunked policy should
execute, segment by segment. To our knowledge no prior work labels states with
measured open-loop divergence, separates open-loop from closed-loop stability
per segment, or uses such labels to schedule the execution horizon.
